\subsection{A post-repair unit obstruction after deeper splitting}

\begin{statement}{A.106}{Proved}
{An explicit degree-nine representative forces a cubic coefficient update}
Put
\[
 \begin{gathered}
 K=\Q(\alpha),\qquad \alpha^9=3,\qquad L=K(\zeta_9),\\
 \tau(\alpha)=\zeta_9\alpha,\qquad\tau(\zeta_9)=\zeta_9,
 \qquad q=321193.
 \end{gathered}
 \tag{A.106.1}
\]
The explicit elements $\beta_{\rm pre},\beta\in K^\times$
defined in (A.106.6) below both have rational norm $-2$
and the canonical two-adic divisor of A.101.
The prime $q$ splits completely in the actual field
\[
 \Q\bigl(\mu_{216},3^{1/27},(-2)^{1/72},
           (\tau^i\beta_{\rm pre})^{1/3}:0\le i<9\bigr).
 \tag{A.106.2}
\]
Thus this prime satisfies the deeper-root, full boundary,
and prior-representative splitting conditions of A.105,
with $n=2,t=3,r=1,M=72$.

For every integer $e$ prime to three and every
$Q,R\in\Q^\times$, put
\[
          E_Q=K(\gamma),\qquad \gamma^3=Q^e\alpha.
 \tag{A.106.3}
\]
The element $\beta R^e$ is a norm from the etale
$K\otimes\Q_q$-algebra $E_Q\otimes\Q_q$ precisely
when $Q\in(\Q_q^\times)^3$.
In particular every global norm preimage of $\beta R^e$
under (A.106.3) requires this local cube condition.
The same necessary condition holds for the norm map
on three-primary completions.
It excludes every update with $3\nmid v_q(Q)$,
for arbitrary $R$, including all updates which allocate
$q$ to the first coefficient at exact three-adic
valuation depth two when $e=8$ and $M=72$.
\end{statement}
\begin{proof}
We give a completely specified polynomial construction
of the two representatives, then compute the local norm
equations without dropping their unit coordinates.

\emph{The interpolation polynomials.}
All of the following integers are fixed:
\[
 \begin{gathered}
 g=5,\quad \theta_0=2827,\quad a_0=146470,
       \quad z_0=100461,\quad \mathcal U=\{1,2,4,5,7,8\},\\
 (v_0,\ldots,v_8)=(0,1,2,2,2,1,1,2,0).
 \end{gathered}
 \tag{A.106.4}
\]
The certificate checks that $q$ is prime, $g$ generates
$\F_q^\times$, $q\equiv1\pmod{216}$,
$\theta_0^{27}=3$, $\theta_0^3=a_0$,
$a_0^9=3$, and that $z_0$ has order nine in $\F_q$.
It also checks $(-2)^{(q-1)/72}=1$ there.
Let $a_i\in\{1,\ldots,q-1\}$ represent $a_0z_0^i$
for $0\le i<9$.
Let $\widetilde a_i\in\{0,\ldots,q^2-1\}$ be its
unique Hensel lift satisfying $\widetilde a_i^9=3$
modulo $q^2$.
For $u\in\mathcal U$, define $\widetilde z_u$ in
the same way from $z_0^u$, with ninth power one.
These definitions involve unique integers; explicitly,
the lift of $x^m=d$ is
$x+q\bigl((d-x^m)/q\bigr)(mx^{m-1})^{-1}$,
where the correction coefficient is reduced modulo $q$.

In $(\Z/q^2\Z)[X,Z]$, form the polynomials
\[
 \begin{split}
 P_i(X)&=\prod_{j\ne i}
                \frac{X-\widetilde a_j}
                     {\widetilde a_i-\widetilde a_j},\\
 J_1(Z)&=\prod_{u\in\mathcal U\setminus\{1\}}
                \frac{Z-\widetilde z_u}
                     {\widetilde z_1-\widetilde z_u}.
 \end{split}
 \tag{A.106.5}
\]
Every denominator is a unit in this coefficient ring.
Let $P_{i,1}\in\Z[X,Z]$ be the unique polynomial of
degree at most eight in $X$ and five in $Z$ whose
coefficients belong to $\{0,\ldots,q^2-1\}$ and whose
reduction is $P_iJ_1$.
This fixes every coefficient, rather than choosing an
unspecified weak-approximation element.
Define in $L$ and then in $K$
\[
 \begin{split}
 h_{\rm pre}&=1+\sum_{i=0}^8(5^{v_i}-1)P_{i,1}(\alpha,\zeta_9),\\
 h_{\rm bad}&=1+(q-1)P_{1,1}(\alpha,\zeta_9)
                         +4P_{4,1}(\alpha,\zeta_9),\\
 u(h)&=N_{L/K}\bigl(\tau(h)/h\bigr),\\
 \beta_{\rm pre}&=(1-\alpha)u(h_{\rm pre}),
       \qquad \beta=\beta_{\rm pre}u(h_{\rm bad}).
 \end{split}
 \tag{A.106.6}
\]
Both $h$ are integral and congruent to one modulo two
in $\Z[\alpha,\zeta_9]$, because all the displayed
coefficients $5^{v_i}-1,q-1,4$ are even.
They are therefore nonzero and are units at every
prime over two.
The formula in (A.106.6) defines actual elements.

Transitivity of field norms and the automorphism
$\tau\in\operatorname{Gal}(L/\Q)$ give the exact identity
\[
 N_{K/\Q}u(h)
      =\frac{N_{L/\Q}\tau(h)}{N_{L/\Q}h}=1.
 \tag{A.106.7}
\]
Also $N_{K/\Q}(1-\alpha)=1-3=-2$.
This proves both rational norms claimed in the statement.
Each $u(h)$ is a unit above two, so both modifications
retain the divisor of $1-\alpha$ there.
The ninth-power root algebra of $3$ over $\Q_2$ is
unramified and has exactly one degree-one factor.
Its root in $\Q_2$ is three modulo eight, giving
$v_2(1-\alpha)=1$ at that factor.
At every other factor the residue of the root differs
from one, so $1-\alpha$ is a unit.
This is precisely the canonical divisor of A.101.

\emph{The exact branch values and the prior splitting.}
Label the $54$ embeddings of $L$ into $\Q_q$ by
$(i,u)\in\Z/9\Z\times\mathcal U$, with residues
$\alpha\mapsto a_i$ and $\zeta_9\mapsto z_0^u$.
They are distinct and exhaustive: the root algebra has
degree $9\varphi(9)=54$, as in A.77.
Their reductions modulo $q^2$ are the fixed Hensel
lifts in (A.106.5).
The polynomials $P_{i,1}$ have value one at $(i,1)$
and zero at all other such embeddings, modulo $q^2$.
At $(i,1)$ the respective $h$ values are consequently
$5^{v_i}$ and $q$ for $i=1$, $5$ for $i=4$, and one
for the other $i$ in the second case.
At every $(i,u)$ with $u\ne1$, both $h$ have value one
modulo $q^2$.

Write cubic characters additively in $\F_3$, choosing
the generator for which the character of $g=5$ is one.
The characters of $1-a_i$ are exactly
\[
       (B_0,\ldots,B_8)=(2,2,0,0,1,0,2,2,0).
 \tag{A.106.8}
\]
The given tuple $(v_i)$ satisfies
$B_i+v_{i+1}-v_i=0$ in $\F_3$ for every cyclic index.
Under the embedding $(i,u)$, $\tau h$ is evaluated
at $(i+u,u)$.
The relative norm in (A.106.6) therefore has valuation
and unit-character vectors obtained by the actual map
\[
 (d_{i,u})_{i,u}\longmapsto
          \left(\sum_{u\in\mathcal U}
                      (d_{i+u,u}-d_{i,u})\right)_{i=0}^8.
 \tag{A.106.9}
\]
The same formula holds separately for normalized unit
characters; at the unique valuation-one $h$ value,
the unit residue is one because that value is $q$
modulo $q^2$.
For $h_{\rm pre}$ all valuations are zero and
(A.106.8)--(A.106.9) cancel every character of
$1-\alpha$.
Hence $\beta_{\rm pre}$ is a cubic unit at all nine
branches of $K$ above $q$.
Together with $q\equiv1\pmod{216}$, the root
$\theta_0$ and the $72$nd-power test for $-2$,
Hensel lifting proves complete splitting in (A.106.2).
In particular every old root $\alpha_i$ is a cube
in $\Q_q$: it has a cube root among the $27$th roots
of three.

Applying (A.106.9) to $h_{\rm bad}$ gives the following
actual valuations and cubic unit characters of $\beta$:
\[
 \begin{split}
 (t_0,\ldots,t_8)&=(1,-1,0,0,0,0,0,0,0),\\
 (C_0,\ldots,C_8)&=(0,0,0,1,2,0,0,0,0),
       \qquad C_i=[q^{-t_i}\beta_i].
 \end{split}
 \tag{A.106.10}
\]
Here $\beta_i$ denotes the image at branch $i$.
Thus all valuations sum to zero and all unit characters
sum to zero, as also follows from its rational norm.

\emph{The necessary and sufficient local norm criterion.}
Write $Q=q^kQ_0$ and $R=q^lR_0$, with unit parts
$Q_0,R_0\in\Z_q^\times$, and put
$x=[Q_0]$, $y=[R_0]$ in $\F_3$.
At branch $i$, the radicand $Q^e\alpha_i$ has
valuation $ek$ and unit character $ex$; the target
$\beta_iR^e$ has valuation $t_i+el$ and unit character
$C_i+ey$.
The complete tame norm calculation (A.103.8) gives
the nine exact equations
\[
            (t_i+el)ex-ek(C_i+ey)=0
                         \quad\hbox{in }\F_3\quad(0\le i<9).
 \tag{A.106.11}
\]
They include split and unramified cubic algebras as
well as the totally ramified case.
Subtract the equations at $i=3$ and $i=4$.
Their equal valuations cancel the first terms, while
$C_3-C_4=-1$ is nonzero.
Since $e$ is a three-adic unit, this forces
$k=0$ modulo three.
With that condition, subtract the equations at $i=0$
and $i=2$; their valuation difference is one.
It follows that $x=0$.
Conversely, $k=x=0$ makes every equation zero, and
then every $Q^e\alpha_i$ is already a cube in $\Q_q$.
The local cubic algebras are split and their norms
are onto.
Because $q\ne3$, a $q$-adic unit is a cube precisely
when its residue is a cube, by Hensel lifting on the
principal units.
Thus $k=x=0$ is exactly $Q\in(\Q_q^\times)^3$.
This proves the local equivalence for every $R$,
without a restriction on its valuation.

Globally (A.106.3) has degree three: the unique
prime of $K$ above three gives radicand valuation
$1+9e\,v_3(Q)$, which is prime to three.
Every global norm localizes to all nine local norms,
so the proved local condition is necessary globally.
Every failed equation in (A.106.11) takes values
in a finite group of exponent three and factors
through the three-primary local completion.
It consequently gives the same obstruction to the
completed global norm map.
\end{proof}

The actual completed norm-fibre map under discussion is
\[
 \{z\in\widehat{E_Q^\times}_3:
          N_{E_Q/\Q}z=-2R^{9e}\}
       \xrightarrow{\ N_{E_Q/K}\ }
 \{w\in\widehat{K^\times}_3:
          N_{K/\Q}w=-2R^{9e}\}.
 \tag{A.106.12}
\]
The target contains $\beta R^e$.
For $e=8$ the corresponding coefficient update is
$a'=3Q^{72}$, $b'=-2R^{72}$, and
$\gamma^{27}=a'$.
An allocation with $3\nmid v_q(Q)$ gives
$v_3(v_q(a'))=2$, but (A.106.11) excludes the
specified target from the displayed norm image.
The ordinary rational pair $(3,-2)$ itself has the
compatible family $1-3^{1/3^j}$ with exact rational
norm $-2$.
This explicit family and (A.106.12) specify the
scope of the representative obstruction.
The local cube condition proved here is not asserted
to suffice for a global preimage.

The exact certificate
\path{code/norm_deeper_unit_certificate.py} and its
retained output check all the residue and interpolation
identities, the two pairs of $54$-dimensional exact
norm determinants, and all $162$ choices of
$e,k,l,x,y$ modulo three.
Exactly $18$ choices satisfy all the equations,
precisely those with $k=x=0$.
The file \path{code/norm_deeper_unit_polynomials.json}
contains every integer coefficient of both $h$ and
their exact rational norms, including the Hensel
lifts used in their definitions.
